\textbf{Phần 1. Trắc nghiệm lựa chọn}
\setcounterpageref{ex}{0}

\begin{ex} %[2D7V33]
Số tròn chục lớn nhất có hai chữ số là
\choice
{\True $90$}
{$99$}
{$80$}
{$100$}
\loigiai{
Số tròn chục lớn nhất có hai chữ số là 90.
\textbf{Đáp án của câu hỏi là A}
}
\end{ex}

\begin{ex} %[2D7V33]
Tìm số nhỏ nhất có hai chữ số mà tổng hai chữ số của nó bằng $10$.
\choice
{$28$}
{$37$}
{\True $19$}
{$18$}
\loigiai{
Để nhỏ nhất, chữ số hàng chục nhỏ nhất là 1, chữ số hàng đơn vị là $10 - 1 = 9$. Số đó là 19.
\textbf{Đáp án của câu hỏi là C}
}
\end{ex}

\begin{ex} %[2D7V33]
Trong số $74$, chữ số $7$ có giá trị lớn hơn chữ số $4$ bao nhiêu đơn vị?
\choice
{$63$ đơn vị}
{\True $66$ đơn vị}
{$70$ đơn vị}
{$3$ đơn vị}
\loigiai{
Giá trị chữ số 7 là 70, chữ số 4 là 4. Hiệu: $70 - 4 = 66$.
\textbf{Đáp án của câu hỏi là B}
}
\end{ex}

\begin{ex} %[2D7V33]
Cho các số: $35, 82, 53, 28, 91$. Số lớn thứ hai trong các số đã cho là
\choice
{$91$}
{$53$}
{$35$}
{\True $82$}
\loigiai{
Sắp xếp giảm dần: $91 > 82 > 53 > 35 > 28$. Số lớn thứ hai là 82.
\textbf{Đáp án của câu hỏi là D}
}
\end{ex}

\begin{ex} %[2D7V33]
Số liền sau của số lớn nhất có hai chữ số khác nhau là
\choice
{\True $99$}
{$98$}
{$100$}
{$97$}
\loigiai{
Số lớn nhất có 2 chữ số khác nhau là 98. Số liền sau là $98 + 1 = 99$.
\textbf{Đáp án của câu hỏi là A}
}
\end{ex}

\begin{ex} %[2D7V33]
Có bao nhiêu số có hai chữ số mà chữ số hàng chục bé hơn chữ số hàng đơn vị $7$ đơn vị?
\choice
{$1$ số}
{$3$ số}
{\True $2$ số}
{$4$ số}
\loigiai{
Các số đó là: 18 ($8 - 1 = 7$) và 29 ($9 - 2 = 7$). Có tất cả 2 số.
\textbf{Đáp án của câu hỏi là C}
}
\end{ex}

\begin{ex} %[2D7V33]
Tìm $x$ biết $x$ là số tròn chục và $45 < x < 65$. Giá trị của $x$ là
\choice
{$50$}
{\True $50\text{ hoặc } 60$}
{$60$}
{$55$}
\loigiai{
Các số tròn chục thỏa mãn là 50 và 60.
\textbf{Đáp án của câu hỏi là B}
}
\end{ex}

\begin{ex} %[2D7V33]
Từ hai chữ số $4$ và $9$, có thể viết được bao nhiêu số có hai chữ số?
\choice
{$2$ số}
{$3$ số}
{$6$ số}
{\True $4$ số}
\loigiai{
Các số là: 44, 49, 94, 99. Có tất cả 4 số.
\textbf{Đáp án của câu hỏi là D}
}
\end{ex}

\begin{ex} %[2D7V33]
Hiệu giữa số chẵn nhỏ nhất có hai chữ số và số lẻ nhỏ nhất có một chữ số là
\choice
{\True $9$}
{$8$}
{$10$}
{$7$}
\loigiai{
Số chẵn nhỏ nhất có 2 chữ số là 10. Số lẻ nhỏ nhất có 1 chữ số là 1. Hiệu: $10 - 1 = 9$.
\textbf{Đáp án của câu hỏi là A}
}
\end{ex}

\begin{ex} %[2D7V33]
Cho dãy số: $5, 10, 15, 20, \dots$ Số hạng thứ $8$ của dãy số là
\choice
{$35$}
{$45$}
{\True $40$}
{$30$}
\loigiai{
Dãy số là các số đếm thêm 5. Số hạng thứ 8 là: $5 \times 8 = 40$.
\textbf{Đáp án của câu hỏi là C}
}
\end{ex}

\begin{ex} %[2D7V33]
Một lớp học có số học sinh nhiều hơn $30$ bạn nhưng ít hơn $35$ bạn. Nếu xếp thành các hàng $4$ bạn thì vừa hết. Số học sinh của lớp là
\choice
{$31$ bạn}
{\True $32$ bạn}
{$33$ bạn}
{$34$ bạn}
\loigiai{
Số chia hết cho 4 nằm giữa 30 và 35 là 32 bạn.
\textbf{Đáp án của câu hỏi là B}
}
\end{ex}

\begin{ex} %[2D7V33]
Điền dấu thích hợp: $50 + 40 \dots 100 - 15$.
\choice
{$<$}
{$=$}
{Không so sánh được}
{\True $>$}
\loigiai{
$50 + 40 = 90$; $100 - 15 = 85$. Vì $90 > 85$ nên điền dấu $>$. 
\textbf{Đáp án của câu hỏi là D}
}
\end{ex}

\textbf{Phần 2. Trắc nghiệm đúng sai}
\setcounterpageref{ex}{0}

\begin{ex} %[2D7V33]
Xét tính đúng sai của các nhận định số học ôn tập học kì 1:
\choiceTF
{\True Số bé nhất có hai chữ số là 10.}
{Số 100 là số lớn nhất có hai chữ số.}
{\True Hai số tự nhiên liên tiếp luôn gồm một số chẵn và một số lẻ.}
{\True Trong số 85, chữ số 8 thuộc hàng chục và có giá trị là 80.}
\loigiai{
\begin{itemize}
    \item Ý a) Đúng: Đúng.
    \item Ý b) Sai: 100 là số có 3 chữ số.
    \item Ý c) Đúng: Đúng.
    \item Ý d) Đúng: Đúng.
\end{itemize}
\begin{itemize}
    \item a) \textbf{Đúng}.
    \item b) \textbf{Sai}.
    \item c) \textbf{Đúng}.
    \item d) \textbf{Đúng}.
\end{itemize}
}
\end{ex}

\begin{ex} %[2D7V33]
Cho ba chữ số 2, 5, 7:
\choiceTF
{\True Lập được tất cả 6 số có hai chữ số khác nhau từ 3 chữ số trên.}
{\True Số lớn nhất lập được là 75.}
{Có 3 số chẵn lập được từ 3 chữ số trên.}
{\True Tổng của số bé nhất (25) và số lớn nhất (75) bằng 100.}
\loigiai{
\begin{itemize}
    \item Ý a) Đúng: Các số là 25, 27, 52, 57, 72, 75.
    \item Ý b) Đúng: Trong các số khác nhau, 75 là lớn nhất.
    \item Ý c) Sai: Chỉ có 2 số có tận cùng là 2: 52 và 72.
    \item Ý d) Đúng: $25 + 75 = 100$. Đúng.
\end{itemize}
\begin{itemize}
    \item a) \textbf{Đúng}.
    \item b) \textbf{Đúng}.
    \item c) \textbf{Sai}.
    \item d) \textbf{Đúng}.
\end{itemize}
}
\end{ex}

\textbf{Phần 3. Trả lời ngắn}
\setcounterpageref{ex}{0}

\begin{ex} %[2D7V33]
Tìm số lớn nhất có hai chữ số mà tổng hai chữ số bằng 11.
\shortans{92}
\loigiai{
\textbf{Đáp số:} $92$\\
Chọn chữ số hàng chục lớn nhất là 9, hàng đơn vị là $11 - 9 = 2$. Số đó là 92.
}
\end{ex}

\begin{ex} %[2D7V33]
Có bao nhiêu số tròn chục có hai chữ số lớn hơn 30?
\shortans{6}
\loigiai{
\textbf{Đáp số:} $6$\\
Các số: 40, 50, 60, 70, 80, 90 (có 6 số).
}
\end{ex}

\begin{ex} %[2D7V33]
Tìm số liền trước của số nhỏ nhất có hai chữ số giống nhau.
\shortans{10}
\loigiai{
\textbf{Đáp số:} $10$\\
Số nhỏ nhất có 2 chữ số giống nhau là 11. Số liền trước là 10.
}
\end{ex}

\begin{ex} %[2D7V33]
Tính tổng các chữ số của số lớn nhất có hai chữ số khác nhau.
\shortans{17}
\loigiai{
\textbf{Đáp số:} $17$\\
Số lớn nhất có 2 chữ số khác nhau là 98. Tổng: $9 + 8 = 17$.
}
\end{ex}

\begin{ex} %[2D7V33]
Từ số 20 đến số 60 có tất cả bao nhiêu số tự nhiên liên tiếp?
\shortans{41}
\loigiai{
\textbf{Đáp số:} $41$\\
$(60 - 20) + 1 = 41$ số.
}
\end{ex}

\begin{ex} %[2D7V33]
Tìm một số biết rằng lấy số đó trừ đi 25 thì được số tròn chục nhỏ nhất có hai chữ số.
\shortans{35}
\loigiai{
\textbf{Đáp số:} $35$\\
Số tròn chục nhỏ nhất là 10. Số đó là: $10 + 25 = 35$.
}
\end{ex}
